\section{Conclusion and Limitations}
\label{sec:conclusion}
We presented \magevl, a lightweight 4B-parameter streaming vision--language model centered on its custom-designed ViT-encoder, \magevit. Built upon the principle of codec-aligned sparsity, \magevit is trained from scratch as a codec-agnostic visual front-end that dynamically extracts motion- and entropy-rich patches across traditional (e.g., HEVC) and neural codecs (e.g., DCVC). Integrated with a causal language decoder and an event-driven System 1 gate, \magevl provides a unified architecture for offline multimodal reasoning, fine-grained spatial grounding, and proactive streaming interaction.

Empirically, \magevl-4B matches flagship baselines on static image understanding while achieving uniform improvements across video QA, temporal grounding, and 2D/3D spatial reasoning benchmarks, accompanied by up to $3.5\times$ wall-clock inference speedups over dense frame sampling. Furthermore, our systematic empirical studies yield seven foundational findings, covering visual pre-training data efficiency, resolution scaling, codec-driven system acceleration, VideoQA SFT redundancy, motion-spatial synergy, AI-driven data pipeline optimization, and Zero-Vision SFT for multimodal RL.



\paragraph{Limitations and Future Work.}
Despite its strong visual and temporal performance, \magevl currently exhibits limitations in complex agentic workflows and mathematical reasoning. These gaps primarily stem from a shortage of high-quality text data in our current training mixture, as well as the omission of RL post-training. In future iterations, we plan to adopt a Zero-Vision SFT alongside RL to bridge these agentic and mathematical capability gaps, while extending \magevl to native audio-visual stream fusion and embodied edge applications.


